\documentclass[11pt]{article}
\usepackage[a4paper,margin=1in]{geometry}
\usepackage{fontspec}
\usepackage[boldfont=false]{xeCJK}
\setCJKmainfont{FandolSong-Regular.otf}
\usepackage{amsmath}
\usepackage{amssymb}
\usepackage{array}
\usepackage{booktabs}
\usepackage{graphicx}
\usepackage{adjustbox}
\usepackage{enumitem}
\setlist[itemize]{topsep=2pt,partopsep=0pt,parsep=0pt,itemsep=2pt,leftmargin=1.4em}
\setlist[enumerate]{topsep=2pt,partopsep=0pt,parsep=0pt,itemsep=2pt,leftmargin=1.6em}
\usepackage{parskip}
\usepackage{fvextra}
\fvset{breaklines=true,breakanywhere=true,fontsize=\small}
\setlength{\parindent}{0pt}

\begin{document}

\begin{center}
  {\LARGE\bfseries Finance and accounting (A Level)}\\[0.35em]
  {\large A Level Business}
\end{center}
\vspace{0.6em}

\section*{Financial statements}

\textbf{Financial statements} 财务报表 are the official records that show how a business is doing with money. You study two of them: the income statement and the statement of financial position. By law, companies must publish these each year.

\section*{The income statement}

The \textbf{income statement} 利润表 shows the \textbf{revenue} 营业收入 and costs of a business over a period (usually a year), and the profit left at each stage.

\[
\text{gross profit} = \text{revenue} - \text{cost of sales}
\]

\begin{itemize}
\item \textbf{cost of sales} 销售成本 is the direct cost of the goods that were sold.

\item \textbf{gross profit} 毛利润 is what is left after taking cost of sales from revenue.

\item after taking off \textbf{expenses} 费用 (such as rent, wages and marketing), you get the \textbf{operating profit} 营业利润.

\item after interest and tax, what is left is the \textbf{net profit} 净利润 (the profit for the year).

\end{itemize}

\section*{The statement of financial position}

The \textbf{statement of financial position} 财务状况表 (also called the \textbf{balance sheet} 资产负债表) is a photo of what the business owns and owes on one day.

\begin{itemize}
\item \textbf{non-current assets} 非流动资产 — items kept for more than a year, such as buildings and machines.

\item \textbf{current assets} 流动资产 — items that become cash within a year, such as inventory and money owed by customers.

\item \textbf{current liabilities} 流动负债 — debts due within a year.

\item \textbf{non-current liabilities} 非流动负债 — debts due after more than a year, such as a long loan.

\item \textbf{equity} 所有者权益 — the money the owners have put in, plus profit kept in the business.

\end{itemize}

The two sides always balance: what the firm owns is funded by what it owes plus the owners' equity.

\section*{Depreciation}

\textbf{Depreciation} 折旧 is the way the cost of a non-current asset is spread over the years it is used, instead of all in the year it was bought. For example, a \$10,000 machine used for five years may lose \$2,000 of value each year. Depreciation lowers the asset's value on the statement of financial position and counts as an expense on the income statement, so profit is not overstated.

\section*{How stakeholders use financial statements}

Different \textbf{stakeholders} 利益相关者 read these statements for different reasons:

\begin{itemize}
\item owners and investors check the profit and the return on their money.

\item lenders check whether the firm can repay its loans.

\item managers use them to make decisions and set budgets.

\item the government checks the tax due; suppliers check the firm can pay.

\end{itemize}

\section*{Ratio analysis}

A single number means little on its own. \textbf{Ratio analysis} 比率分析 links two figures from the statements to judge performance, and lets you compare one year with another, or one firm with another.

\section*{Profitability ratios}

\textbf{Profitability} 盈利能力 ratios show how good the firm is at turning sales and capital into profit.

\[
\text{gross profit margin} = \frac{\text{gross profit}}{\text{revenue}} \times 100\%
\]

\[
\text{operating profit margin} = \frac{\text{operating profit}}{\text{revenue}} \times 100\%
\]

A higher \textbf{gross profit margin} 毛利率 or \textbf{operating profit margin} 营业利润率 means more profit is kept from each sale. The key overall measure is \textbf{return on capital employed} 资本回报率 (ROCE), which compares profit with the money invested:

\[
\text{ROCE} = \frac{\text{operating profit}}{\text{capital employed}} \times 100\%
\]

Here \textbf{capital employed} 所用资本 is the total long-term money in the business. A higher ROCE means the money is being used well.

\section*{Liquidity ratios}

\textbf{Liquidity} 流动性 means how easily a firm can pay its short-term debts. Two ratios test it.

\[
\text{current ratio} = \frac{\text{current assets}}{\text{current liabilities}}
\]

\[
\text{acid test ratio} = \frac{\text{current assets} - \text{inventory}}{\text{current liabilities}}
\]

The \textbf{current ratio} 流动比率 of about 1.5–2 is usually safe. The \textbf{acid test ratio} 速动比率 is stricter because it removes inventory, which can be slow to sell. Too low is risky; too high may mean cash is sitting idle.

\section*{Efficiency and gearing ratios}

\textbf{Efficiency} 效率 ratios show how well the firm uses its resources — for example, \textbf{inventory turnover} 存货周转率 measures how many times a year the firm sells and replaces its stock. Faster turnover usually ties up less cash.

\textbf{Gearing} 杠杆比率 shows how much of the firm's long-term money comes from borrowing rather than from owners:

\[
\text{gearing} = \frac{\text{non-current liabilities}}{\text{capital employed}} \times 100\%
\]

High gearing means heavy borrowing — risky if interest rates rise, but it can boost returns when times are good.

\section*{Limitations of ratio analysis}

Ratios are useful but limited. They use past data, which may not predict the future. They ignore non-money factors like staff morale and brand strength. And a fair comparison needs firms of similar size in the same industry.

\section*{Investment appraisal}

\textbf{Investment appraisal} 投资评估 is the study of whether a big spending project (such as a new machine or factory) is worth it. Three methods are used.

\section*{Payback, ARR and NPV}

\begin{itemize}
\item the \textbf{payback period} 回收期 is the time it takes for the project's cash inflows to repay the initial cost. A shorter payback is safer.

\item the \textbf{accounting rate of return} 会计回报率 (ARR) shows the average yearly profit as a percentage of the money invested:

\end{itemize}

\[
\text{ARR} = \frac{\text{average annual profit}}{\text{initial investment}} \times 100\%
\]

\begin{itemize}
\item the \textbf{net present value} 净现值 (NPV) recognises that money in the future is worth less than money today. It uses \textbf{discounting} 折现 to shrink future cash flows back to today's value, then takes away the initial cost. A positive NPV means the project adds value.

\end{itemize}

Each method has strengths: payback is simple and focuses on risk; ARR shows profitability; NPV is the most complete but needs a chosen discount rate.

\section*{Finance and accounting strategy}

At the top level, managers use financial data and budgets to choose between strategies. Before a big decision, they ask: can we afford it, how will it affect profit, cash flow and gearing, and what is the likely return? Financial information turns a risky guess into a reasoned choice — but numbers are only part of the picture, and must be weighed with the market, the staff and the firm's objectives.


\end{document}
